% Capítulo 8 — El álgebra de las transformaciones: productos, determinante, autovectores y formas lineales
\chapter{El álgebra de las transformaciones: productos, determinante, autovectores y formas lineales}\label{cap:08}\label{sec:bloque-algebraico}

\section{La idea}

Primero el mapa: cada operación del álgebra se lee como una pregunta geométrica. Lo armamos con los resúmenes del profesor y dejamos las cuentas para la sección siguiente.

\subsection{Del álgebra a la geometría: qué pregunta cada operación}\label{sec:algebra-geometria}

\textbf{Desarrollo.} El hilo de la materia es siempre el mismo: una operación algebraica,
en apariencia abstracta, resulta ser exactamente la pregunta geométrica que necesitábamos.
Conviene tener el mapa antes de las cuentas.

\begin{center}
\begin{tabularx}{\textwidth}{>{\bfseries}p{34mm} X X}
\toprule
Operación & Qué calcula & Qué pregunta geométrica responde\\
\midrule
Producto interno \(u\cdot v\) & un número & ¿son ortogonales? ¿cuánto se proyecta uno
sobre otro? En homogéneas: ¿el punto está en la recta?\\
Producto vectorial \(u\times v\) & un vector normal a ambos & ¿cuál es la recta que une dos
puntos? ¿cuál es el punto donde se cortan dos rectas? ¿qué área orientada forman?\\
Determinante \(\det[u\ v\ w]\) & un número & ¿están alineados tres puntos? ¿son
linealmente dependientes tres vectores? ¿qué volumen orientado forman?\\
Producto exterior \(u\wedge v\) & un bivector (área orientada) & lo mismo que el
vectorial, pero en cualquier dimensión\\
Núcleo de \(A\) & un subespacio & ¿qué vectores aplasta \(A\) a cero? Si hay alguno,
\(\det A=0\)\\
Autovectores de \(A\) & vectores con \(Av=\lambda v\) & ¿qué direcciones deja quietas
\(A\)? En homogéneas: ¿qué puntos deja fijos la transformación?\\
\bottomrule
\end{tabularx}
\end{center}

\begin{idea}[La cadena]
incidencia \(\to\) unión de puntos \(\to\) alineación \(\to\) determinante \(\to\)
núcleo \(\to\) autovectores \(\to\) puntos fijos. Cada eslabón se deduce del anterior;
es lo que desarrollan las secciones siguientes de este capítulo.
\end{idea}

\begin{idea}[En palabras del profesor]
``Traemos a la mesa a nuestro `detector de colinealidad': el producto cuña, y lo usamos
primero para construir un polinomio cuyas raíces son los autovalores, y luego para hallar
los autovectores como el subespacio nulo de unas matrices que armamos'' (resumen de clase 7, cursada 2025).
\end{idea}

\section{Construcciones, cuentas y ejemplos}

Hacemos la cuenta de cada operación con su figura: producto interno, producto vectorial, determinante, núcleo, autovalores y formas lineales. Lo dictado fija el orden; los resúmenes del profesor y nuestro desarrollo completan los pasos.

\textbf{Dictado} (clase del 18 de septiembre). El objetivo declarado por el docente: que la
parte geométrica y la parte algebraica ``se conecten'', y entender por qué el determinante
tiene las propiedades ``mágicas'' que tiene. Lo que sigue reconstruye ese razonamiento en el
orden en que se dio.

\subsection{Producto interno: incidencia}\label{sec:interno}

{\small\vspace{1mm}\linkUnidad{02_producto_interno}{Simulación interactiva: Unidad 4 — Producto interno: proyección, longitud, ángulo}\par}


\textbf{Desarrollo.} El producto interno sirve para dos cosas: proyectar, y decidir si dos
vectores son ortogonales (\(u\cdot v=0\)). En coordenadas homogéneas aparece un tercer uso.
Un punto se representa por una terna \(p=(x,y,w)^T\) y una recta \(ax+by+c=0\) por otra
terna \(\ell=(a,b,c)^T\). Son objetos distintos (el docente los dibuja con símbolos
distintos), pero ambos son vectores de \(\R^3\), y
\[
  \boxed{\ \ell\cdot p=ax+by+cw=0\quad\Longleftrightarrow\quad p\ \text{está en}\ \ell\ }.
\]
La ecuación cartesiana de la recta es literalmente ``el producto interno con la normal es
cero''. Ni \(\ell\) ni \(p\) cambian de significado si se los multiplica por un número, y la
condición tampoco.

\subsection{Producto vectorial: unir dos puntos, cortar dos rectas}\label{sec:vectorial}

\textbf{Desarrollo.} Queremos una operación que, dados dos puntos \(p,q\), devuelva la recta
\(\ell=p\times q\) que los une. ¿Qué tiene que cumplir? Que \(p\) esté en esa recta y \(q\)
también:
\[
  (p\times q)\cdot p=0,\qquad (p\times q)\cdot q=0 .
\]
Por lo tanto la operación debe producir un vector \emph{ortogonal a sus dos argumentos}. Y
más: si ponemos dos veces el mismo punto no hay recta que determinar, así que \(p\times p=0\).
Estas exigencias, junto con la linealidad en cada argumento, \emph{obligan} a que la
operación sea, a menos de una constante, el producto vectorial de \(\R^3\) (en homogéneas
la constante no importa):
\[
  p\times q=\begin{pmatrix}p_2q_3-p_3q_2\\ p_3q_1-p_1q_3\\ p_1q_2-p_2q_1\end{pmatrix}.
\]

\begin{resultado}[Las dos propiedades que lo definen]
\begin{itemize}
  \item \textbf{Linealidad} en cada argumento: \((2p)\times q=2\,(p\times q)\),
  \((p+p')\times q=p\times q+p'\times q\).
  \item \textbf{Antisimetría}: \(p\times q=-\,q\times p\). Es equivalente a \(p\times p=0\):
  el único vector igual a su opuesto es el nulo.
\end{itemize}
Son exactamente las propiedades del \emph{área orientada} de un paralelogramo (duplicar un
lado duplica el área; cambiar el orden de los lados cambia la orientación; dos lados
paralelos no encierran área). Por eso \(\lVert p\times q\rVert\) es el área del paralelogramo
de lados \(p,q\).
\end{resultado}

Por simetría de la ecuación de incidencia, la misma operación resuelve el problema dual:
\emph{el punto donde se cortan dos rectas es} \(\ell\times m\). Dos rectas paralelas
\((a,b,c)\) y \((a,b,c')\) dan \(\ell\times m=(c'-c)\,(b,-a,0)^T\): un punto ideal, su
dirección común.

\subsection{Tres puntos alineados: el determinante}\label{sec:determinante}

{\small\vspace{1mm}\linkUnidad{07a_cuna_determinante}{Simulación interactiva: Unidad 9 — Cuña y determinante: el detector de colinealidad}\par}


\textbf{Desarrollo.} ¿Está el punto \(r\) en la recta que une \(p\) y \(q\)? Se forma la
recta y se pregunta por la incidencia:
\[
  (p\times q)\cdot r=0 .
\]
Su anulación es simétrica en los tres puntos (si \(r\) está en \(pq\), entonces \(p\) está
en \(qr\) y \(q\) en \(rp\)); el valor, en cambio, es antisimétrico: se anula si dos de
ellos coinciden, es lineal en cada uno y cambia de signo al intercambiar dos. Es, término a término, el determinante:
\[
  (p\times q)\cdot r=\det\begin{pmatrix}p_1&p_2&p_3\\q_1&q_2&q_3\\r_1&r_2&r_3\end{pmatrix}.
\]

\begin{resultado}[El determinante es la única operación posible]
Una operación de tres vectores de \(\R^3\) que sea lineal en cada argumento y que se anule
cuando dos argumentos coinciden es, a menos de una constante, el determinante. Vale cero
exactamente cuando los tres vectores son linealmente dependientes (coplanares como vectores;
alineados como puntos homogéneos) y su valor absoluto es el volumen del paralelepípedo que
forman. Intercambiar dos filas cambia el signo: es la antisimetría otra vez.
\end{resultado}

\paragraph{Verificación numérica hecha en clase.} Con \(p=(1,2,3)\), \(q=(4,5,-1)\),
\(r=(0,0,1)\):
\[
  p\times q=\begin{pmatrix}2\cdot(-1)-3\cdot5\\ 3\cdot4-1\cdot(-1)\\ 1\cdot5-2\cdot4\end{pmatrix}
  =\begin{pmatrix}-17\\13\\-3\end{pmatrix},\qquad
  (p\times q)\cdot r=-3,
\]
y el determinante de la matriz de filas \(p,q,r\), desarrollado por la tercera fila, es
\(1\cdot(1\cdot5-2\cdot4)=-3\). Coinciden, como debía ser: desarrollar por una fila
\emph{es} hacer un producto vectorial con las otras dos y luego un producto interno.

\subsection{Núcleo de una matriz y determinante nulo}\label{sec:nucleo}

{\small\vspace{1mm}\linkUnidad{07b_nucleo_autovectores}{Simulación interactiva: Unidad 10 — Núcleo y autovectores: las direcciones que no giran}\par}


\textbf{Desarrollo.} El \emph{núcleo} de \(A\) es el conjunto de vectores \(v\) con
\(Av=0\). Si se piensa \(Av\) fila por fila, cada coordenada del resultado es el producto
interno de una fila de \(A\) con \(v\). Que las tres den cero significa que \(v\) es
ortogonal a las tres filas: las tres filas están en el plano de normal \(v\), son linealmente
dependientes, y por lo tanto
\[
  \boxed{\ \text{núcleo}(A)\neq\{0\}\quad\Longleftrightarrow\quad\det A=0\ }.
\]
Es la propiedad del determinante que se usa para encontrar autovalores.

\subsection{Autovalores y autovectores}\label{sec:autovalores}

{\small\vspace{1mm}\linkUnidad{07c_autovalores_3d_puntos_fijos}{Simulación interactiva: Unidad 11 — Autovalores en 3D y puntos fijos proyectivos}\par}


\textbf{Dictado.} \(v\neq0\) es \emph{autovector} de \(A\) con \emph{autovalor}
\(\lambda\) si \(Av=\lambda v\): la transformación lo deja sobre su propia recta. En
coordenadas homogéneas eso es exactamente ``la transformación deja fijo el punto que \(v\)
representa''.

\paragraph{El pase de magia.} \(Av=\lambda v\) equivale a \(Av-\lambda Iv=0\), es decir
\[
  (A-\lambda I)\,v=0 :
\]
\(v\) está en el núcleo de \(A-\lambda I\). Para que ese núcleo no sea trivial,
\[
  \boxed{\ \det(A-\lambda I)=0\ }\qquad\text{(polinomio característico)} .
\]
Para \(A\) de \(3\times3\) es un polinomio de grado \(3\) en \(\lambda\).

\begin{resultado}[Cuántos autovalores reales]
\begin{itemize}
  \item Un polinomio real de grado \(3\) tiene tres raíces, todas reales, o una real y dos
  complejas conjugadas. Nunca tres complejas: las raíces complejas de un polinomio real
  vienen de a pares conjugados.
  \item A cada autovalor real le corresponde un autovector real. A un autovalor complejo no
  le corresponde ningún vector de \(\R^3\).
  \item Una matriz de rotación de \(\R^3\) tiene siempre el autovalor \(1\): su autovector es
  el eje. Los otros dos son \(e^{\pm i\theta}\), con \(\theta\) el ángulo de giro; son reales
  sólo para \(\theta=0^\circ\) (identidad: \(1,1,1\)) y \(\theta=180^\circ\) (\(1,-1,-1\)),
  y complejos conjugados en cualquier otro caso.
\end{itemize}
\end{resultado}

\begin{idea}[Método del producto cruz para el autovector]
Encontrado \(\lambda\), las filas de \(A-\lambda I\) son todas ortogonales al autovector.
Entonces el autovector es el producto vectorial de dos de esas filas:
\[
  v=(\text{fila}_i)\times(\text{fila}_j) .
\]
Si esas dos filas son paralelas y el producto da cero, se toma otro par. Si \emph{todos} los
pares dan cero, \(\operatorname{rango}(A-\lambda I)\le1\): el autoespacio tiene dimensión
\(\ge2\) (multiplicidad \emph{geométrica} \(\ge2\)) y hay al menos todo un plano de
autovectores (en homogéneas: toda una recta de puntos fijos). La vuelta no vale: una raíz
doble del polinomio (multiplicidad \emph{algebraica} \(2\)) no garantiza que los pares se
anulen; ver la matriz 3a de la tabla de testigos más abajo. No hace falta resolver el sistema: es la
misma operación de ``unir dos rectas en su punto común'' de la sección~\ref{sec:vectorial}.
\end{idea}

\subsubsection{Ejemplo 1: reflexión respecto del plano \texorpdfstring{$xy$}{xy}}

\(A=\operatorname{diag}(1,1,-1)\) conserva \(x,y\) e invierte \(z\).
\[
  \det(A-\lambda I)=(1-\lambda)^2(-1-\lambda)=0
  \quad\Longrightarrow\quad\lambda=1\ (\text{doble}),\ \lambda=-1 .
\]
Para \(\lambda=-1\): \(A+I=\operatorname{diag}(2,2,0)\), y el producto cruz de las dos
primeras filas es \((2,0,0)\times(0,2,0)=(0,0,4)\): el eje \(z\), como se esperaba, con
\(A(0,0,4)^T=(0,0,-4)^T\). Para \(\lambda=1\): \(A-I=\operatorname{diag}(0,0,-2)\), y
cualquier par de filas da producto cruz nulo. Es la raíz doble: todo el plano \(xy\) es de
autovectores, porque la reflexión lo deja quieto punto a punto.

\subsubsection{Ejemplo 2: eje de una rotación compuesta}

Se compone una rotación de \(90^\circ\) alrededor de \(y\) seguida de una de \(90^\circ\)
alrededor de \(z\):
\[
  R=R_z(90^\circ)\,R_y(90^\circ)=
  \begin{pmatrix}0&-1&0\\1&0&0\\0&0&1\end{pmatrix}
  \begin{pmatrix}0&0&1\\0&1&0\\-1&0&0\end{pmatrix}
  =\begin{pmatrix}0&-1&0\\0&0&1\\-1&0&0\end{pmatrix}.
\]
Es una rotación, así que \emph{sin calcular el polinomio} se sabe que \(\lambda=1\) es
autovalor y que su autovector es el eje. Filas de \(R-I\): \((-1,-1,0)\), \((0,-1,1)\),
\((-1,0,-1)\). Producto cruz de las dos primeras:
\[
  (-1,-1,0)\times(0,-1,1)=(-1,\,1,\,1),
  \qquad
  R\begin{pmatrix}-1\\1\\1\end{pmatrix}=\begin{pmatrix}-1\\1\\1\end{pmatrix}\ \checkmark
\]
El eje es la dirección \((-1,1,1)\). Los otros autovalores son \(-\frac12\pm\frac{\sqrt3}2 i\):
un giro de \(120^\circ\) alrededor de ese eje. Con ángulos genéricos (por ejemplo
\(R_z(30^\circ)R_y(50^\circ)\)) la receta es idéntica y se hace en planilla o en NumPy: el
eje resulta \((-0.2263,\,0.8446,\,0.4853)\), y \(Rv=v\) se verifica multiplicando.

\begin{atencion}[Sobre los ``números lindos'']
Los ejercicios de álgebra están diseñados para que den enteros. Una matriz genérica, como
la de una rotación de \(37^\circ\) o la homografía de un pentágono regular, da raíces
irracionales y autovalores con muchos decimales. Eso no es un error: es lo normal. Por eso
el método debe funcionar en planilla o en código, no sólo a mano.
\end{atencion}

\subsection{Formas lineales: la recta como máquina de medir}\label{sec:formas-lineales}

{\small\vspace{1mm}\linkUnidad{06_formas_lineales}{Simulación interactiva: Unidad 8 — Formas lineales y dualidad}\par}
{\small\vspace{1mm}\linkAnimacion{01_dualidad_gradiente}{Animación: dualidad punto–gradiente}\par}


\textbf{Desarrollo} (a partir de los resúmenes de clase 4 y 9 del profesor). Hasta acá los
vectores fueron flechas: cosas que se suman y se estiran. Hay otro objeto, igual de simple,
que en los libros aparece tarde y que el profesor prefiere presentar apenas definido el
espacio: la \emph{forma lineal}, una función que toma un vector y devuelve un número,
respetando sumas y escalas.

\paragraph{Primero en \(\R\).} Una forma lineal en la recta es multiplicar por un número:
\(\varphi(x)=2x\). Es una regla graduada: los puntos donde \(\varphi\) vale
\(0,1,2,3,\ldots\) están equiespaciados sobre la recta, a distancia \(\tfrac12\). No hay
nada más que saber. Al pasar al plano la regla sigue siendo la misma, pero ahora hay
direcciones que la regla no mide.

\paragraph{Un ejemplo con números chicos.} Sea
\[
  \varphi(x,y)=2x+3y .
\]
Con \(v=(1,0)\) da \(2\); con \((0,1)\) da \(3\); con \((1,1)\) da \(5\); con \((3,-2)\)
da \(0\). Es lineal: \(\varphi(v+w)=\varphi(v)+\varphi(w)\) y
\(\varphi(\lambda v)=\lambda\,\varphi(v)\). Y basta conocer \(\varphi\) en los dos vectores
de la base para conocerla en todos: \(\varphi(x,y)=x\,\varphi(1,0)+y\,\varphi(0,1)\).

\paragraph{Hecho geométrico: una regla graduada tirada en el plano.} Los puntos donde
\(\varphi\) vale lo mismo forman una recta, \(2x+3y=c\), y al variar \(c\) se obtiene una
familia de rectas paralelas equiespaciadas. La forma lineal \emph{es} esa familia: para
medir un punto, se mira qué recta de la familia pasa por él. La recta \(c=0\) pasa por el
origen y contiene la dirección \((3,-2)\), que es lo que la forma ``no ve'' (su núcleo).
El vector \((2,3)\) es perpendicular a todas las rectas de nivel y apunta hacia donde
\(\varphi\) crece.

\begin{center}
\begin{tikzpicture}[scale=.78]
  \begin{scope}
    \clip (-3.2,-3.2) rectangle (4.2,3.2);
    \foreach \c in {-6,-4,-2,0,2,4,6,8}
      \draw[teal,line width=.7pt] (-3.2,{(\c+6.4)/3})--(4.2,{(\c-8.4)/3});
    \draw[teal,line width=1.2pt] (-3.2,{6.4/3})--(4.2,{-8.4/3});
    \draw[axisline] (-3.2,0)--(4.2,0);
    \draw[axisline] (0,-3.2)--(0,3.2);
    \draw[rust,line width=1pt,-{Latex[length=2.4mm]}] (0,0)--(1,1.5)
      node[right,font=\small,labelback] {\((2,3)\)};
    \draw[tealdark,line width=1pt,-{Latex[length=2.4mm]}] (0,0)--(1.5,-1)
      node[below right,font=\small,labelback] {\((3,-2)\): \(\varphi=0\)};
    \fill[ink] (1,1) circle (2pt) node[above left,font=\small,labelback] {\((1,1)\): \(\varphi=5\)};
  \end{scope}
  \foreach \c/\y in {8/-0.13,6/-0.8,4/-1.47,2/-2.13}
    \node[right,font=\scriptsize,text=teal] at (4.2,\y) {\(c=\c\)};
  \foreach \c/\y in {0/2.13,-2/1.47,-4/0.8,-6/0.13}
    \node[left,font=\scriptsize,text=teal] at (-3.2,\y) {\(c=\c\)};
\end{tikzpicture}
\end{center}

\paragraph{Las selectoras y el espacio de las formas.} Las formas más básicas son las que
seleccionan una coordenada: \(\mathsf{x}(x,y)=x\) y \(\mathsf{y}(x,y)=y\). Cualquier otra
es combinación de ellas: \(\varphi=2\,\mathsf{x}+3\,\mathsf{y}\). Las formas lineales se
suman y se multiplican por escalares, y el resultado sigue siendo una forma lineal: forman
un espacio vectorial, el \emph{dual} del plano. En la base canónica una forma se describe
con dos números, igual que un vector, pero conviene escribirla como \emph{fila}:
\[
  \varphi(v)=\begin{pmatrix}2&3\end{pmatrix}\begin{pmatrix}x\\y\end{pmatrix}
  =(2,3)\cdot v .
\]
La fila por la columna es el producto interno con el vector \((2,3)\): la forma se
``disfraza'' de vector gracias al producto interno. En coordenadas homogéneas esto ya se
usó: una recta \(\ell=(a,b,c)\) es una forma lineal, un punto \(p\) es un vector, y
\(\ell\cdot p=0\) es la incidencia de la sección~\ref{sec:interno}.

\paragraph{La matriz como pila de formas.} Una transformación lineal del plano al plano
devuelve dos números por cada vector; cada uno de ellos es una forma lineal. Con
\[
  A=\begin{pmatrix}2&3\\1&-1\end{pmatrix},\qquad
  A\begin{pmatrix}1\\2\end{pmatrix}=\begin{pmatrix}2\cdot1+3\cdot2\\1\cdot1-1\cdot2\end{pmatrix}
  =\begin{pmatrix}8\\-1\end{pmatrix},
\]
la primera fila mide con la regla \(2x+3y\) y la segunda con la regla \(x-y\).
Geométricamente, aplicar \(A\) es leer dos reglas graduadas a la vez. El núcleo de \(A\)
es el conjunto de puntos donde las dos reglas marcan cero: la intersección de las rectas
\(2x+3y=0\) y \(x-y=0\), que es sólo el origen porque las rectas se cortan
(\(\det A=-5\neq0\)). Si en cambio \(B=\begin{pmatrix}2&3\\4&6\end{pmatrix}\), las dos
filas son la misma regla escalada, sus rectas de nivel cero coinciden, el núcleo es toda
la recta \(2x+3y=0\) y \(\det B=0\): es la sección~\ref{sec:nucleo} vista desde las
filas. Y el método del producto cruz para autovectores (sección~\ref{sec:autovalores})
busca justamente el vector que todas las reglas de \(A-\lambda I\) marcan en cero.

\paragraph{Cómo se transforma una forma cuando se transforma el vector.} Sea
\(A=\begin{pmatrix}1&1\\0&2\end{pmatrix}\) y \(v=(1,1)\). Entonces \(Av=(2,2)\) y
\(\varphi(Av)=2\cdot2+3\cdot2=10\). ¿Qué regla, aplicada al \(v\) original, da ese mismo
\(10\)? Se calcula la fila por la matriz:
\[
  (2,3)\,A=\bigl(2\cdot1+3\cdot0,\ 2\cdot1+3\cdot2\bigr)=(2,8),
  \qquad (2,8)\cdot(1,1)=10\ \checkmark
\]
El vector se multiplica por \(A\) a la izquierda (columna); la forma se multiplica por
\(A\) a la derecha (fila). Es lo que el profesor resume como \(\varphi(A\vec v)=(\varphi A)(\vec v)\):
``una forma lineal se transforma del lado contrario que un vector''.

Si lo que se quiere es \emph{arrastrar la regla junto con la figura}, es decir, encontrar
la forma \(\psi\) cuyas rectas de nivel son las imágenes por \(A\) de las rectas de nivel
de \(\varphi\), hay que deshacer \(A\): \(\psi=\varphi A^{-1}\). Con
\(A^{-1}=\begin{pmatrix}1&-\tfrac12\\0&\tfrac12\end{pmatrix}\),
\[
  \psi=(2,3)\,A^{-1}=\bigl(2,\ -1+\tfrac32\bigr)=\bigl(2,\tfrac12\bigr).
\]
Verificación: el punto \((3,-2)\) está en la recta \(\varphi=0\); su imagen es
\(A(3,-2)=(1,-4)\), y \(\psi(1,-4)=2-2=0\): sigue en la recta de nivel cero de la regla
arrastrada. El punto \((3,0)\), con \(\varphi=6\), va a \((3,0)\) y \(\psi(3,0)=6\).
Hecho geométrico: al deformar el plano las rectas de nivel se deforman con él, siguen
siendo rectas paralelas equiespaciadas (es una afinidad), pero la separación y la
inclinación cambian, y por eso los coeficientes de la forma cambian ``al revés'' que los
de los vectores.

\paragraph{El gradiente es una forma lineal.} Sea \(f(x,y)=x^2+y\) en el punto \((1,1)\),
donde \(f=2\). Para un desplazamiento chico \((h,k)\),
\[
  f(1+h,1+k)=2+2h+k+h^2\approx 2+\underbrace{(2h+k)}_{\text{forma lineal}} .
\]
La derivada en \((1,1)\) es la regla \(2h+k\): la mejor aproximación lineal de \(f\) cerca
del punto. Su recta de nivel cero, \(2(x-1)+(y-1)=0\), es la tangente a la curva de nivel
\(x^2+y=2\) en \((1,1)\). Que ``el gradiente es perpendicular a la curva de nivel'' es la
misma figura de arriba: el vector \((2,1)\) es perpendicular a la familia de rectas de la
forma. Ir a lo chico (la forma lineal en un punto) para volver a lo grande (la superficie
completa) es lo que hace el capítulo~\ref{cap:11} con el plano tangente y la métrica.

\begin{idea}[Aplicación en Inteligencia Artificial: Support Vector Machines (SVM) como formas lineales]
Desarrollo complementario, desde la especialización en IA (CEIA); no es de la cátedra. Un
\emph{clasificador lineal} separa dos lotes de puntos del plano (dos clases, \(y_i=\pm1\))
con una regla \(\omega(x)=w_1x+w_2y+b\): la clase es el signo de \(\omega\). Es exactamente
una forma lineal, una regla graduada tirada sobre la nube: los puntos viven en \(V\), el
clasificador vive en \(V^{*}\), y en coordenadas homogéneas \((x,y,1)\) la traslación \(b\)
entra en la misma fila, \(\omega=(w_1,w_2,b)\), igual que una recta en la
sección~\ref{sec:interno}. La máquina de vectores soporte elige, entre todas las reglas
que separan, la que deja más \emph{margen}: pide \(\omega(x_i)\ge+1\) para una clase y
\(\le-1\) para la otra, y entre las que cumplen toma la de \(\|w\|\) mínima. Las tres
rectas de nivel \(-1\), \(0\), \(+1\) son el borde de un lado, la frontera y el borde del
otro; los puntos que tocan los bordes son los \emph{vectores soporte}. El ancho de la
franja entre las rectas \(-1\) y \(+1\) es \(2/\|w\|\), porque la separación entre
curvas de nivel consecutivas de una forma es \(1/\|w\|\) (clase del 30/9: ``cuanto más
largo el vector de la forma, más apretadas las curvas''), y esa cuenta necesita el
producto interno: sin métrica hay regla, pero no hay ancho. Ejemplo con enteros: clase
\(+1\) en \((1,1)\) y \((2,0)\), clase \(-1\) en \((-1,-1)\) y \((-2,0)\). La frontera
de margen máximo es \(x+y=0\) con \(w=(\tfrac12,\tfrac12)\), \(b=0\): los cuatro puntos
dan \(\omega=\pm1\) exacto (los cuatro son vectores soporte) y el ancho es
\(2/\|w\|=2\sqrt2\); la recta vertical \(x=0\) también separa, pero con ancho \(2\). Lo
verificamos resolviendo el problema de minimización con \texttt{scipy.optimize}. La unidad
\texttt{06\_formas\_lineales} (apéndice~\ref{ap:notebooks}) lo dibuja con el helper
\texttt{svm\_margen\_duro\_2d} de \file{cb104.py}: mover un punto y ver cómo gira la regla.
\end{idea}

\begin{resultado}[Dualidad en pocas palabras]
\begin{itemize}
  \item Un vector es una columna; una forma lineal es una fila. Fila por columna es un
  número: la medición.
  \item Hecho geométrico: vector = flecha; forma = familia de rectas paralelas
  equiespaciadas (una regla). En \(\R^3\), planos paralelos.
  \item Una matriz es una pila de formas (sus filas) o una hilera de vectores (sus
  columnas). Ambas lecturas son útiles: las filas dan el núcleo, las columnas dan la imagen.
  \item Al transformar el plano con \(A\), los vectores van con \(A\) y las formas con
  \(A^{-1}\) (o con \(A\) del otro lado, si se pregunta ``qué regla mide lo mismo'').
  \item En homogéneas, los puntos son vectores y las rectas son formas; cambiarlos de rol
  es la dualidad de la geometría proyectiva.
\end{itemize}
\end{resultado}

\section{En profundidad}

Extendemos lo anterior con el producto exterior en cualquier dimensión y con el cambio de base y la diagonalización, que se dictaron el 30 de septiembre. Lo que excede lo dictado queda marcado como anticipado.

\subsection{Producto exterior: la misma idea en cualquier dimensión}\label{sec:exterior}

\textbf{Desarrollo.} El producto vectorial existe sólo en \(\R^3\). La operación que
generaliza la idea a cualquier dimensión es el \emph{producto exterior} de Grassmann,
\(u\wedge v\), definido únicamente por linealidad y antisimetría (\(u\wedge u=0\)). Su
resultado no es un vector sino un \emph{bivector}: un elemento de área orientada.

Escribiendo \(u=ae_1+be_2+ce_3\) y \(v=de_1+ee_2+fe_3\) y distribuyendo, los términos
\(e_i\wedge e_i\) se anulan y \(e_2\wedge e_1=-e_1\wedge e_2\), de modo que
\[
  u\wedge v=(ae-bd)\,e_1\wedge e_2+(af-cd)\,e_1\wedge e_3+(bf-ce)\,e_2\wedge e_3 .
\]
Son los mismos tres números del producto vectorial; sólo cambia el signo del término del
medio, porque la regla de la mano derecha asigna \(e_1\times e_3=-e_2\). Los ``planos
elementales'' \(e_1\wedge e_2\), \(e_1\wedge e_3\), \(e_2\wedge e_3\) son la base de las
áreas orientadas de \(\R^3\).

Operando con un tercer vector \(w=ge_1+he_2+ie_3\) sólo sobreviven los términos que completan
las tres direcciones, y reordenando (cada intercambio de vecinos cambia el signo) todo cae
sobre un único elemento de volumen:
\[
  (u\wedge v)\wedge w=\det[u\ v\ w]\;e_1\wedge e_2\wedge e_3 .
\]
De ahí la regla de los signos \(+,-,+\) del desarrollo por menores.

\begin{idea}[Contar subespacios]
En \(\R^3\) hay \(3\) direcciones de vectores, \(3\) planos elementales y \(1\) volumen; por
eso el resultado final es un número. En \(\R^4\) hay \(4\) direcciones, \(6\) planos, \(4\)
volúmenes tridimensionales y \(1\) hipervolumen. La operación es mecánica en cualquier
dimensión: da una fórmula cerrada para un vector ortogonal a un hiperplano de \(\R^n\), útil
para autovectores en dimensión baja cuando \(\operatorname{rango}(A-\lambda I)=n-1\). Para
matrices grandes los autovectores se calculan con otros algoritmos (QR, Lanczos).
\end{idea}

\paragraph{Ortogonal es lo que maximiza el volumen.} Entre todos los vectores de norma
\(1\), el que forma mayor volumen con un plano dado es el perpendicular al plano; los otros
producen un prisma ``inclinado'' de menor altura. El docente adelantó que de esta idea sale
la regla para encontrar el vector ortogonal a un subespacio en cualquier dimensión.

\subsection{Cambio de base y diagonalización}\label{sec:cambio-base}

{\small\vspace{1mm}\linkUnidad{08_cambio_de_base}{Simulación interactiva: Unidad 12 — Cambio de base y diagonalización}\par}
{\small\vspace{1mm}\linkAnimacion{03_cambio_de_base_viaje}{Animación: el viaje \(A\,T_{\mathrm{diag}}\,A^{-1}\)}\par}


\begin{atencion}[Anticipado]
El cambio de base y la diagonalización se dictaron el 30 de septiembre de 2026 (clase 8 del
esqueleto) y se desarrollan en el párrafo ``La matriz de cambio de base y el despeje de la
diagonalización''. Lo que sigue sobre potencias \(A^t\) y trayectorias no se dictó en 2C2026: se
desarrolla antes de que la cátedra lo dicte, a partir del resumen 2025 del profesor y de las
fichas de clase. Pendiente de validar con lo que efectivamente se dicte.
\end{atencion}

\paragraph{Hecho geométrico: la flecha no se movió, cambió la grilla.} Un mojón tiene
coordenadas \((3,4)\) en la grilla cuadrada de siempre. Si en vez de los ejes canónicos usamos
las direcciones \(b_1=(1,1)\) y \(b_2=(1,-1)\), el \emph{mismo} punto se escribe
\[
  \begin{pmatrix}3\\4\end{pmatrix}
  =\tfrac72\begin{pmatrix}1\\1\end{pmatrix}-\tfrac12\begin{pmatrix}1\\-1\end{pmatrix},
  \qquad\text{coordenadas nuevas }(\tfrac72,-\tfrac12).
\]
El mojón está donde estaba; cambió la regla con que lo medimos. Escalera: en \(\R\) cambiar
de base es cambiar de unidad (de metros a centímetros el número se multiplica por 100 y la
longitud no cambia); en \(\R^2\) la grilla puede ser oblicua y con pasos distintos; en
\(\R^n\) es una matriz \(n\times n\) invertible. El profesor observó en 2025 que ``no parece
haber un conjunto de vectores privilegiado a priori'': lo único que distingue a la base
canónica es que con ella el producto interno es la cuenta de siempre.

\begin{resultado}[La matriz de cambio de base]
Si \(P\) tiene en sus columnas los vectores de la base nueva escritos en la canónica,
\[
  [v]_{\text{can}}=P\,[v]_{\text{nueva}},\qquad [v]_{\text{nueva}}=P^{-1}[v]_{\text{can}} .
\]
Con \(P=\begin{pmatrix}1&1\\1&-1\end{pmatrix}\), \(P^{-1}=\tfrac12\begin{pmatrix}1&1\\1&-1\end{pmatrix}\)
y \(P^{-1}(3,4)^T=(\tfrac72,-\tfrac12)^T\), como arriba. Los tres tipos de objeto cambian
de base distinto, y conviene escribir siempre el sentido (acá, de la canónica a la nueva):
los vectores como \([v]_{\text{nueva}}=P^{-1}[v]_{\text{can}}\); las formas lineales, que
son filas, multiplicando por \(P\) del lado derecho, \([\varphi]_{\text{nueva}}=[\varphi]_{\text{can}}\,P\)
(sección~\ref{sec:formas-lineales}), para que \(\varphi(v)\) dé lo mismo en las dos bases;
y las transformaciones con
\[
  A_{\text{nueva}}=P^{-1}\,A\,P ,
\]
que se lee de derecha a izquierda: pasamos las coordenadas nuevas a canónicas, aplicamos \(A\)
como siempre, volvemos a las nuevas. Todo se convierte a la vez para que ninguna relación
cambie: cada punto sigue en su recta y \(A\) sigue llevando cada punto adonde lo llevaba.
\end{resultado}

\paragraph{El ejercicio del 30 de septiembre, resuelto.} Consigna del profesor (en hoja
grande, dibujado): formas lineales \((1,1)\) y \((-1,2)\), vector \((3,4)\), transformación
\(T=\begin{pmatrix}1{,}2&0{,}5\\-0{,}1&1{,}3\end{pmatrix}\), elegida ``no muy distinta de
la identidad''. Dibujar las curvas de nivel de las dos formas (dos mazos de rectas
paralelas que juntos arman una grilla), ubicar el vector en la grilla, transformar el
vector con \(T\) y las formas con \(T^{-1}\), redibujar y verificar que el vector
transformado cae en la misma intersección. Antes de transformar, \((1,1)\cdot(3,4)=7\) y
\((-1,2)\cdot(3,4)=5\): el vector está en la séptima curva de una y la quinta de la otra.
Inversa por adjunta sobre determinante, \(\det T=1{,}56+0{,}05=1{,}61\):
\[
  T^{-1}=\frac1{1{,}61}\begin{pmatrix}1{,}3&-0{,}5\\0{,}1&1{,}2\end{pmatrix},\qquad
  T\begin{pmatrix}3\\4\end{pmatrix}=\begin{pmatrix}5{,}6\\4{,}9\end{pmatrix},
\]
\[
  (1,1)\,T^{-1}=\tfrac1{1{,}61}(1{,}4,\ 0{,}7)=(0{,}870,\ 0{,}435),\qquad
  (-1,2)\,T^{-1}=\tfrac1{1{,}61}(-1{,}1,\ 2{,}9)=(-0{,}683,\ 1{,}801).
\]
Verificación: \(0{,}870\cdot5{,}6+0{,}435\cdot4{,}9=7{,}0\) y
\(-0{,}683\cdot5{,}6+1{,}801\cdot4{,}9=5{,}0\). El vector transformado sigue en la séptima
curva de la primera forma y en la quinta de la segunda. En clase la forma \((1,1)\)
quedó ``\((0{,}86,\ 0{,}43)\), se acostó un poco'' (redondeo) y la \((-1,2)\) ``un poco más
vertical''. Hecho geométrico: la grilla se deformó con el plano y el punto se movió con el
plano; medido con la grilla deformada, el punto está donde estaba. ``No hacemos nada,
convertimos todo para que nada cambie'' (profesor, 30/9/2026). En planilla: dos columnas
para la forma, dos para el vector, \texttt{SUMAPRODUCTO}.

\paragraph{Diagonalizar es elegir la grilla de los autovectores.} Sea
\(A=\begin{pmatrix}2&1\\1&2\end{pmatrix}\). Polinomio característico
\((2-\lambda)^2-1=0\), autovalores \(3\) y \(1\), autovectores \((1,1)\) y \((1,-1)\)
(perpendiculares porque \(A\) es simétrica). Con esos autovectores como columnas de \(P\),
\[
  P^{-1}AP=\tfrac12\begin{pmatrix}1&1\\1&-1\end{pmatrix}
  \begin{pmatrix}3&1\\3&-1\end{pmatrix}
  =\begin{pmatrix}3&0\\0&1\end{pmatrix}=D .
\]
Hecho geométrico: en la grilla de los autovectores la transformación es dos estiramientos
independientes, por \(3\) a lo largo de \(b_1\) y por \(1\) a lo largo de \(b_2\); la grilla
``se estira o se comprime pero no se rompe''. En la grilla común la misma transformación
no tiene lectura fácil. Es lo que el resumen 2025 describe como ``aplicar dos afinidades
axiales sucesivas''.

\paragraph{La matriz de cambio de base y el despeje de la diagonalización.} En este
párrafo usamos los nombres de la clase: \(A\) es la matriz que tiene los autovectores en
columnas (la \(P\) del recuadro de arriba, con otro nombre) y \(T\) es la transformación
(la \(A\) del párrafo anterior).

\emph{El diagrama de conmutatividad.} Dualidad y cambio de base se juntan en una sola
cadena, escrita en orden de aplicación (de izquierda a derecha):
\[
  (p)\;\to\;(A)\;\to\;[\,A^{-1}\to T_{\mathrm{diag}}\to A\,]\;\to\;(A^{-1})\;\to\;(\ell) .
\]
En un extremo entra un punto \(p\) y en el otro una recta \(\ell\), los dos escritos en la
base de autovectores; lo que se mide es la incidencia. Con la convención de abajo (\(A\)
lleva de la base de autovectores a la canónica), el punto pasa a la canónica como \(A\,p\)
y la forma como \(\ell\,A^{-1}\) (multiplicando del lado derecho). El recuadro del medio
recibe y devuelve coordenadas canónicas: adentro actúa \(T_{\mathrm{diag}}\), en la base
de autovectores, y el recuadro entero, leído como producto (lo que se aplica último va a la
izquierda), es \(A\,T_{\mathrm{diag}}\,A^{-1}=T_{\mathrm{can}}\), la misma transformación
escrita en la base canónica. La cadena completa es
\(\ell A^{-1}\,(A\,T_{\mathrm{diag}}\,A^{-1})\,A\,p=\ell\,T_{\mathrm{diag}}\,p\): las
conversiones se absorben de a pares (\(A^{-1}A=I\)) y la incidencia no cambia; el punto
transformado sigue cayendo en la recta transformada. Despejar
\(T_{\mathrm{diag}}=A^{-1}\,T_{\mathrm{can}}\,A\) es sacar la diagonal del recuadro, como
hacemos más abajo. Todo se convierte a la vez para que nada cambie.

\emph{La matriz con los autovectores en columnas.} Definimos \(A=[\,v_1\mid v_2\,]\), los
autovectores escritos en la base canónica, uno por columna, y la probamos con el vector
\((1,0)^T\):
\[
  A\begin{pmatrix}1\\0\end{pmatrix}
  =\begin{pmatrix}A_{11}&A_{12}\\A_{21}&A_{22}\end{pmatrix}\begin{pmatrix}1\\0\end{pmatrix}
  =\begin{pmatrix}A_{11}\\A_{21}\end{pmatrix}=v_1 .
\]
Entra una unidad del primer autovector escrita en la base de autovectores y sale ese mismo
autovector escrito en la canónica: \(A\) traduce de la base de autovectores a la canónica
y \(A^{-1}\) hace el camino inverso, de la canónica a la de autovectores. El recorrido
completo para un punto cualquiera: el punto \((1,1)=e_x+e_y\), dado en la canónica, pasa
con \(A^{-1}\) a coordenadas \((c_1,c_2)\) en la base de autovectores, es decir
\((1,1)=c_1v_1+c_2v_2\); ahí la transformación es diagonal y devuelve
\((\lambda_1c_1,\ \lambda_2c_2)\); y \(A\) traduce ese resultado de vuelta a la canónica,
\(\lambda_1c_1v_1+\lambda_2c_2v_2\), que es \(T(1,1)\). Con la matriz del párrafo
anterior, \(v_1=(1,1)\) y \(v_2=(1,-1)\), el punto \((1,1)\) es el propio \(v_1\):
\((c_1,c_2)=(1,0)\), la diagonal da \((3,0)\) y \(A\) devuelve \(3v_1=(3,3)\), que
coincide con \(\begin{pmatrix}2&1\\1&2\end{pmatrix}(1,1)^T=(3,3)^T\).

\emph{El despeje.} El mismo recorrido, escrito como producto de matrices, dice que la
transformación en la base canónica es la diagonal conjugada por el cambio de base:
\[
  T_{\mathrm{can}}=A\,T_{\mathrm{diag}}\,A^{-1}
  \quad\Longrightarrow\quad
  T_{\mathrm{diag}}=A^{-1}\,T_{\mathrm{can}}\,A ,
\]
leída de derecha a izquierda: \(A^{-1}\) lleva a la base de autovectores,
\(T_{\mathrm{diag}}\) estira cada eje propio, \(A\) vuelve a la canónica. El despeje se
hace multiplicando por \(A^{-1}\) a la izquierda y por \(A\) a la derecha; los dos
productos \(A^{-1}A=I\) y \(A\,A^{-1}=I\) que aparecen se cancelan y queda
\(T_{\mathrm{diag}}\) sola. Con los nombres de este apunte (\(P\) para la matriz de
autovectores, \(A\) para la transformación) es exactamente el \(D=P^{-1}AP\) del párrafo
anterior: misma matriz, otro nombre.

\paragraph{Potencias y la trayectoria de un punto.} Volvemos a los nombres del apunte:
\(P\) es la matriz de autovectores y \(A\) la transformación. Como \(A=PDP^{-1}\),
\(A^n=PD^nP^{-1}\): aplicar \(n\) veces la transformación es estirar \(n\) veces cada eje
propio. Para \(n=5\),
\[
  A^5=P\begin{pmatrix}243&0\\0&1\end{pmatrix}P^{-1}
  =\tfrac12\begin{pmatrix}244&242\\242&244\end{pmatrix}
  =\begin{pmatrix}122&121\\121&122\end{pmatrix}.
\]
Lectura: \((1,0)=\tfrac12(1,1)+\tfrac12(1,-1)\); la parte en \((1,1)\) se multiplica por
\(243\) y la parte en \((1,-1)\) queda igual: \(\tfrac12(243,243)+\tfrac12(1,-1)=(122,121)\),
la primera columna. Si los autovalores son positivos el exponente puede ser real:
\(A^t=PD^tP^{-1}\) con \(D^t=\operatorname{diag}(3^t,1)\), y el punto \((1,0)\) recorre
\[
  A^t\begin{pmatrix}1\\0\end{pmatrix}=\tfrac12\begin{pmatrix}3^t+1\\3^t-1\end{pmatrix}:
  \quad(1,0)\ \text{en } t=0,\quad(1{,}37,\,0{,}37)\ \text{en } t=\tfrac12,\quad(2,1)\ \text{en } t=1 .
\]
Es la tercera idea del hilo (sección~\ref{sec:hilo}): la transformación se identifica con
la trayectoria que genera, y la potencia real la convierte en película. La misma jugada que
la exponencial de una rotación, ahora con estiramientos en vez de giros. La unidad
\texttt{08\_cambio\_de\_base} (apéndice~\ref{ap:notebooks}) tiene el deslizador en \(t\).

\begin{atencion}[Una convención, tres escrituras en el resumen 2025]
Verificado en el resumen del profesor: la clase 7 escribe
\(A=B_{av\mapsto e}\,D\,B_{e\mapsto av}\) y la clase 10 escribe
\(A=C^{-1}_{e\mapsto av}\,\Lambda\,C_{e\mapsto av}\); en ambas el factor de la derecha,
el que actúa primero sobre un vector en canónica, es el que \emph{convierte de \(e\) a
\(av\)}. La clase 8 escribe \(M_{av}=C_{av\mapsto e}\,M\,C_{e\mapsto av}\), donde con la
misma lectura el factor de la derecha debería ser el que convierte de \(av\) a \(e\)
(un vector en base \(av\) entra primero). Los subíndices de la clase 8 están al revés de
los de las clases 7 y 10. En la clase del 30 de septiembre de 2026 quedó con subíndices
descriptivos, no direccionales:
\(A\,T_{\mathrm{diag}\langle Av\rangle}A^{-1}=T_{\mathrm{no\,diag}\langle\mathrm{can}\rangle}\)
con \(A\) = autovectores en columnas, que es \(P\,D\,P^{-1}\) de arriba. Para no mezclar, en este apunte \(P\) siempre tiene la base
nueva en columnas y se lee ``\(P\) lleva de la nueva a la canónica''.
\end{atencion}

\begin{resultado}[Diagonalizabilidad no tiene nada que ver con invertibilidad: cuatro matrices testigo]
Las dos preguntas de la clase 6, ``¿qué pasa si hay algo en el núcleo?'' y ``¿qué pasa si
un autovalor repetido no alcanza para armar una base?'', son independientes. Invertir
pide que el núcleo sea \(\{0\}\), es decir, que ningún autovalor sea \(0\). Diagonalizar
pide una base de \(\R^2\) formada por autovectores. Las cuatro combinaciones aparecen:
\medskip

\renewcommand{\arraystretch}{1.35}
\noindent\begin{tabular}{@{}p{2.6cm}p{1.1cm}p{2.2cm}p{2.1cm}p{5.6cm}@{}}
\textbf{Matriz} & \(\det\) & \textbf{Autovalores} & \textbf{¿Base de autovectores?} & \textbf{Hecho geométrico}\\
\hline
\textbf{1.}~\(\bigl(\begin{smallmatrix}2&0\\0&1\end{smallmatrix}\bigr)\) & \(2\ne0\) & \(2,\ 1\) & sí, \(e_1,e_2\) &
diagonalizable e invertible: estira el eje \(x\) al doble y deja fijo el \(y\); dos direcciones propias independientes.\\
\textbf{2.}~\(\bigl(\begin{smallmatrix}1&0\\0&0\end{smallmatrix}\bigr)\) & \(0\) & \(\lambda_1=1,\ \lambda_2=0\) & sí, \(e_1,e_2\) &
proyección sobre el eje \(x\): aplasta el eje \(y\) al origen, sin inversa, pero con base completa de autovectores.\\
\textbf{3a.}~\(\bigl(\begin{smallmatrix}1&1\\0&1\end{smallmatrix}\bigr)\) & \(1\) & \(1\) (doble) & no: sólo \(e_1\) &
corte: invertible, con inversa \(\bigl(\begin{smallmatrix}1&-1\\0&1\end{smallmatrix}\bigr)\); sólo el eje \(x\) queda fijo.\\
\textbf{3b.}~\(\bigl(\begin{smallmatrix}0&-1\\1&0\end{smallmatrix}\bigr)\) & \(1\) & \(\pm i\) & no en \(\R\) &
giro de \(90^\circ\): invertible, ninguna recta real queda fija.\\
\textbf{4.}~\(\bigl(\begin{smallmatrix}0&1\\0&0\end{smallmatrix}\bigr)\) & \(0\) & \(0\) (doble) & no: sólo \(e_1\) &
ni diagonalizable ni invertible: manda \(e_2\mapsto e_1\) y \(e_1\mapsto0\).\\
\end{tabular}
\renewcommand{\arraystretch}{1.18}
\medskip

\textbf{Matriz 2.} Tener autovalor \(0\) impide invertir (el núcleo es el eje \(y\),
\(\det=0\)) y no impide diagonalizar: \(e_1\) y \(e_2\) son autovectores con
\(\lambda_1=1\), \(\lambda_2=0\), y en esa base la matriz ya es diagonal.

\textbf{Matriz 3a.} El polinomio característico es \((\lambda-1)^2=0\): la raíz
\(\lambda=1\) tiene multiplicidad algebraica \(2\). Pero \((A-I)v=0\) da \(v_2=0\), de modo
que el autoespacio es sólo el eje \(x\): multiplicidad geométrica \(1\). Falta un segundo
autovector, y sin él no hay base.

\textbf{Matriz 3b.} En el giro el motivo es otro: no hay rectas reales fijas, los
autovalores \(\pm i\) no son reales.

\textbf{Matriz 4.} \(\lambda=0\) doble, pero el núcleo es sólo el eje \(x\): un único
autovector \(e_1\).

El defecto de las matrices 3a, 3b y 4 es exactamente por lo que la sección~\ref{sec:svd}
pasa a la SVD, \(A=U\Sigma V^T\): con una base ortonormal a la entrada (\(V\)) y otra
\emph{distinta} a la salida (\(U\)), toda matriz, incluso un corte o una rectangular, se
lee como diagonal. Es diagonal entre dos grillas distintas, no una semejanza: \(\Sigma\) no
es ``\(A\) en otra base'', los \(\sigma\) no son autovalores (en el corte,
\(\sigma=(\sqrt5\pm1)/2\) y \(\lambda=1\) doble) y no sirve para potencias,
\(A^n\neq U\Sigma^nV^T\). Las dos descomposiciones coinciden sólo si \(A\) es simétrica
semidefinida positiva.
\linkUnidad{09_cambio_de_base_svd}{Simulación interactiva: cambio de base y SVD}
\end{resultado}
